%MB9T5
\loai{Thay đổi bước sóng. Tính khoảng vân}
\begin{vidu}%[3P5-3.5K][Loai1]
Trong thí nghiệm Young về giao thoa ánh sáng với ánh sáng đơn sắc có bước sóng ${{\lambda}_{1}}=540\ nm $ thì thu được hệ vân giao thoa trên màn quan sát có khoảng vân ${i_1}=0,36\ mm $. Khi thay ánh sáng trên bằng ánh sáng đơn sắc có bước sóng ${{\lambda}_{2}}=600\ nm $ thì thu được hệ vân giao thoa trên màn quan sát có khoảng vân
\choice
{${i_2}=0,60\ mm $}
{\True ${i_2}=0,40\ mm $}
{${i_2}=0,50\ mm $}
{${i_2}=0,45\ mm $}
\loigiai{
\begin{itemize}
\item Ta có: $\dfrac{i_1}{i_2}=\dfrac{{{\lambda}_{1}}}{{{\lambda}_{2}}}\Rightarrow {i_2}=\dfrac{{{\lambda}_{2}}{i_1}}{{{\lambda}_{1}}}=\dfrac{600. 0,36}{540}=0,4\ mm $.
\end{itemize}
}
\end{vidu} %\dotlineEX{4}

\loai{Thay đổi bước sóng. Tính tổng hợp}
\begin{vidu}%[3P5-3.5K][Loai3]
Trong thí nghiệm Young về giao thoa ánh sáng, nguồn sáng phát ra ánh sáng đơn sắc có bước sóng ${{\lambda}_{1}} $. Trên màn quan sát, trên đoạn thẳng MN dài 20 mm (MN vuông góc với hệ vân giao thoa) có 10 vân tối, M và N là vị trí của hai vân sáng. Thay ánh sáng trên bằng ánh sáng đơn sắc có bước sóng ${{\lambda}_{2}}=\dfrac{5{{\lambda}_{1}}}{3} $ thì tại M là vị trí của một vân sáng giao thoa, số vân sáng trên đoạn MN lúc này là
\choice
{\True 7}
{5}
{8}
{6}
\loigiai{\begin{itemize}
\item Theo bài ra $MN=10{i_1}=20\Rightarrow {i_1}=2\ mm $ Ta có: $\dfrac{i_1}{i_2}=\dfrac{{{\lambda}_{1}}}{{{\lambda}_{2}}}\Rightarrow {i_2}=\dfrac{{{\lambda}_{2}}{i_1}}{{{\lambda}_{1}}}=\dfrac{5. 2}{3}=\dfrac{10}{3}\ mm $.
\item Ta thấy $\dfrac{MN}{i_2}=6 $ mà tại M là 1 vân giao thoa (vân sáng) vậy có $6 + 1 = 7$ vân sáng.
\end{itemize}
}
\end{vidu} %\dotlineEX{4}
\loai{Thay đổi a}
\begin{vidu}%[3P5-3.5G][Loai4]
Trong thí nghiệm Young về giao thoa ánh sáng, hai khe được chiếu bằng ánh sáng đơn sắc $\lambda $, màn quan sát cách mặt phẳng hai khe một khoảng không đổi D, khoảng cách giữa hai khe có thể thay đổi (nhưng ${S_1} $ và ${S_2} $ luôn cách đều S). Xét điểm M trên màn, lúc đầu là vân sáng bậc 3, nếu lần lượt giảm hoặc tăng khoảng cách ${S_1}{S_2} $ một lượng $\Delta a $ thì tại đó là vân sáng bậc k và bậc 5k. Nếu tăng khoảng cách ${S_1}{S_2} $ thêm $3\Delta a $ thì tại M là
\choice
{vân sáng bậc 7}
{\True vân sáng bậc 9}
{vân sáng bậc 8}
{vân tối thứ 9}
\loigiai{\begin{itemize}
\item $OM=3i=ki'=5ki''\Rightarrow \dfrac{3}{a}=\dfrac{k}{a-\Delta a}=\dfrac{5k}{a+\Delta a}=\dfrac{k'}{a+3\Delta a}\Rightarrow \Delta a=\dfrac{2}{3}a $.
\item $\dfrac{3}{a}=\dfrac{k'}{a+3\Delta a}=\dfrac{k'}{3a}\Rightarrow k'=9 $.
\item Vậy tại M là vân sáng bậc 9.
\end{itemize}
}
\end{vidu} %\dotlineEX{5}
\loai{Dịch màn ra xa. Tìm bước sóng}
\begin{vidu}%[3P5-3.5G][Loai5]
\nguon{ĐH - 2013} Thực hiện thí nghiệm Young về giao thoa với ánh sáng có bước sóng $\lambda$. Khoảng cách giữa hai khe hẹp là 1\ mm. Trên màn quan sát, tại điểm M cách vân trung tâm 4,2\ mm có vân sáng bậc 5. Giữ cố định các điều kiện khác, di chuyển dần màn quan sát dọc theo đường thẳng vuông góc với mặt phẳng chứa hai khe ra xa cho đến khi vân giao thoa tại M chuyển thành vân tối lần thứ hai thì khoảng dịch màn là 0,6 m. Bước sóng $\lambda$ bằng
\choice
{\True $0,6\ \mu m$}
{$0,5\ \mu m$}
{$0,7\ \mu m$}
{$0,4\ \mu m$}
\loigiai{
\begin{itemize}
\item Vị trí điểm M: $x_M=5i=5\dfrac{\lambda D}{a}=4,2.10^{-3}\ m \quad (1)$.
\item Ban đầu, các vân tối tính từ vân trung tâm đến M lần lượt có tọa độ là 0,5i; 1,5i; 2,5i; 3,5i và 4,5i. 
\item Khi dịch màn ra xa 0,6 m, M trở thành vân tối thứ 2 thì:\\ $x_M=3,5i'$ hay $x_M=3,5\dfrac{\lambda \left({D+0,6}\right)}{a}=4,2.10^{-3}\ m\quad (2)$.
\item Từ $(1)$ và $(2)$ tính ra: $D = 1,4\ m$, $\lambda =0,6\ \mu m$.
\end{itemize}
}
\end{vidu} %\dotlineEX{4}
\loai{Dịch màn ra xa. Tính tổng hợp}
\begin{vidu}%[3P5-3.5G][Loai6]
\nguon{QG - 2017} Trong thí nghiệm Young về giao thoa ánh sáng đơn sắc, khoảng cách từ mặt phẳng chứa hai khe đến màn quan sát là 1,5 m. Trên màn quan sát, hai điểm M và N đối xứng qua vân trung tâm có hai vân sáng bậc 4. Dịch màn ra xa hai khe thêm một đoạn 50 cm theo phương vuông góc với mặt phẳng chứa hai khe. So với lúc chưa dịch chuyển màn, số vân sáng trên đoạn MN lúc này giảm đi
\choice
{6 vân}
{7 vân}
{\True 2 vân}
{4 vân}
\loigiai{
\begin{itemize}
\item Số vân sáng trên đoạn MN lúc đầu (ứng với $k_1 = 4$) là: $\Rightarrow N_{s1}=\dfrac{MN}{i}+1=\dfrac{8i}{i}+1=9$ vân sáng.
\item Ta có: $i=\dfrac{\lambda D}{a}\Rightarrow
\begin{cases}
i_{1}=\dfrac{\lambda {D}_{1}}{a}\\
i_{2}=\dfrac{\lambda {D}_{2}}{a}\end{cases}
\Rightarrow \dfrac{{i}_{2}}{{i}_{1}}=\dfrac{{D}_{2}}{{D}_{1}}=\left({\dfrac{2}{1,5}}\right)=\dfrac{4}{3}$.
\item Tại M: $x_{{S}_{1}}=x_{{S}_{2}}\Rightarrow k_{1}i_{1}=k_{2}i_{2}\Rightarrow \dfrac{{k}_{1}}{{k}_{2}}=\dfrac{{i}_{2}}{{i}_{1}}=\dfrac{{D}_{2}}{{D}_{1}}=\dfrac{4}{3}\xrightarrow{{k}_{1}=4}k_{2}=3$.
\item Vậy số vân sáng lúc này là 7 vân, so với lúc đầu giảm đi 2 vân.
\end{itemize}
}
\end{vidu} %\dotlineEX{4}
\loai{Dịch chuyển màn lại gần}
\begin{vidu}%[3P5-3.5G][Loai7]
\nguon{QG - 2016} Trong thí nghiệm Young về giao thoa ánh sáng đơn sắc, khoảng cách hai khe không đổi. Khi khoảng cách từ mặt phẳng chứa hai khe tới màn quan sát là D thì khoảng vân trên màn là 1 mm. Khi khoảng cách từ mặt phẳng chứa hai khe tới màn quan sát lần lượt là $(D -\Delta D)$ và $(D +\Delta D)$ thì khoảng vân trên màn tương ứng là i và 2i. Khi khoảng cách từ mặt phẳng chứa hai khe tới màn quan sát là $(D + 3\Delta D)$ thì khoảng vân trên màn là
\choice
{3 mm}
{3,5 mm}
{\True 2 mm}
{2,5 mm}
\loigiai{\begin{itemize}
\item Do $\lambda$ và D không đổi nên khoảng vân i tỉ lệ với khoảng cách D, ta có hệ:
\begin{align*}
\begin{cases}
\dfrac{\lambda}{a}D=1\ mm\,\,\,\left({1}\right)\\
\dfrac{\lambda}{a}\left({D-\Delta D}\right)=i\,\,\,\left({2}\right)\\
\dfrac{\lambda}{a}\left({D+\Delta D}\right)=2i\,\,\,\left({3}\right)\\
\dfrac{\lambda}{a}\left({D+3\Delta D}\right)=i'\,\,\,\left({4}\right)
\end{cases}
\end{align*}
\item Từ $(2)$ và $(3)$ suy ra $D=3\Delta D$.
\item Từ $(4)$ và $(1)$ suy ra: $i'=\dfrac{D+3\Delta D}{D}=1+\dfrac{3\Delta D}{D}=2\ mm$.
\end{itemize}
}
\end{vidu} %\dotlineEX{5}
\loai{Tổng hợp}
\begin{vidu}%[3P5-3.5K][Loai9]
Trong thí nghiệm Young về giao thoa ánh sáng, hai khe được chiếu bằng ánh sáng đơn sắc có bước sóng ${{\lambda}_{1}}=400\ nm $ thì khoảng vân là ${i_1} $. Nếu tăng khoảng cách giữa màn và mặt phẳng hai khe lên gấp đôi đồng thời thay nguồn sáng phát ánh sáng bước sóng ${{\lambda}_{2}} $ thì khoảng vân là ${i_2}=3{i_1} $. Bước sóng ${{\lambda}_{2}} $ có giá trị
\choice
{\True 0,6 $\mu m $}
{0,5 $\mu m $}
{0,75 $\mu m $}
{0,56 $\mu m $}
\loigiai{
\begin{itemize}
\item Ta có: ${i_2}=3{i_1}\Rightarrow \dfrac{{{\lambda}_{2}}. 2D}{a}=3\dfrac{{{\lambda}_{1}}D}{a}\Rightarrow {{\lambda}_{2}}=\dfrac{3}{2}{{\lambda}_{1}}=\dfrac{3}{2}. 0,4=0,6\ \mu m $.
\end{itemize}
}
\end{vidu} %\dotlineEX{5}
